\documentclass[10pt,a4paper]{amsart}
\usepackage[margin=19mm]{geometry}
\usepackage{amsmath,amssymb,mathtools}
\usepackage[scr=rsfs,cal=euler]{mathalfa}
\usepackage{xfrac}
\usepackage[unicode,hidelinks,bookmarks=false]{hyperref}
\newcommand{\ie}{i.e.\ }
\newcommand{\cf}{cf.\ }
\newcommand{\resp}{respectively\ }
\newcommand{\Subsection}{\subsection}
\providecommand{\nobreakdash}{\nobreak\hbox{-}}
\setlength{\emergencystretch}{2em}
\multlinegap0pt
\usepackage{amssymb}
\usepackage[shortlabels]{enumitem}
\usepackage{float}
\usepackage{mathtools}
\usepackage{tikz}
\newtheorem{thm}{Theorem}
\newcommand{\Cb}{\operatorname{Cb} }
\newcommand{\Pb}{\operatorname{Pb} }
\title{Commuting probability of skew left braces}
\author{Susanta Mondal, Manoj K. Yadav}
\date{Source: arXiv:2603.16771v1 (CC BY 4.0)}
\begin{document}
\maketitle

\noindent\emph{Source and licence.} Commuting probability of skew left braces. Version arXiv:2603.16771v1 (CC BY 4.0), distributed by arXiv under the Creative Commons Attribution 4.0 International licence, http://creativecommons.org/licenses/by/4.0/, as recorded in the arXiv licence metadata for this exact version and shown on the abstract page https://arxiv.org/abs/2603.16771v1. Full licence notice: \texttt{licenses/arxiv-2603.16771-CC-BY-4.0.txt}.\par\smallskip

\noindent\textit{Editorial scope.} Counted unit: the article's thm cbd (source0.tex lines 549--551). The article's complete original proof of it is delivered with it (source0.tex lines 552--585, one \texttt{proof} environment of 34 lines). No mathematics is written, repaired or re-derived: the statement and the proof are the article's own lines, in the article's own notation, and no retained line is substituted or re-flowed.  No further source-local context is needed: every cross-reference of every retained line resolves inside the two delivered spans.  Nothing was omitted from any retained span, so the package carries no omission record.  No retained line contains a citation, so the wrapper carries no bibliography and no bibliography entry.  No retained line contains a figure environment, an image-inclusion command or drawing code, so no image is delivered and the package declares no asset dependency.  Editorial formatting only, in addition: an AMS \texttt{amsart} preamble that restates the article's own declarations and macros used by the retained lines, namely the environments \texttt{thm} on separate counters, together with the standard packages those lines need.  The retained environments print in this document as Theorem 1 in order of appearance, whereas the article prints the same items with the numbering of its own full document.  The complete native source of the article as distributed in the arXiv e-print for this exact version is bundled unchanged as \texttt{sources/196576/source0.tex}.
\begin{thm}\label{Ideal and cbd}
Let $(B,+, \circ)$ be a  skew  brace of finite order and $N$ be an ideal of $B$. Then $\Pb(B)\leqslant \Pb(N)\Pb(B/N)$. In particular, $\Pb(B)\leqslant \Pb(B/N)$ and $\Pb(B) \leqslant \Pb(N)$.
\end{thm}

\begin{proof}
We work with group $(B, \circ)$ and the quotient group $(B, \circ)/(N, \circ)$. For the simplicity of notation we suppress the use of `$\circ$' and write $x \circ N = xN$ for a left coset of $N$ in $(B, \circ)$. For the product of two subgroups $N_1$ and $N_2$  of $(B, \circ)$ we simply write $N_1 N_2$.

 We first prove that $({\Cb_{B}(x)  N)/N} \subset \Cb_{B/N}(x  N)$ for all $x\in B$. Let $y \in \Cb_{B}(x)$.  Since $N$ is an ideal of $B$, we know that 
  $x + N = x \circ N$ for all $x \in B$. Thus by an easy computation we get 
 \begin{eqnarray*}
 (y  N) * (x  N) &=& (y*x)  N  = N,\\
  (x  N) * (y  N) &=& (x*y)  N  = N
  \end{eqnarray*}
  and 
  $$[x N, y  N]^{\circ} = [x, y]^{\circ}  N  = N.$$
Hence $({\Cb_{B}(x)  N)/N} \subset \Cb_{B/N}(x N)$.

Note that $N \cap \Cb_B(x) = \Cb_N(x)$ for all $x \in B$. Thus, as $\Cb_B(x)$ and  $\Cb_N(x)$ are subgroups of $(B, \circ)$, we get
$$(\Cb_B(x)  N)/N \cong \Cb_B(x)/(\Cb_B(x) \cap N) = \Cb_B(x)/\Cb_N(x).$$
Now, using the definition of $\Pb(B)$, we have
\begin{align*}
  |B|^2\Pb(B) &=\sum_{x\in B}|\Cb_{B}(x)| =\sum_{b N\in{B/N}} ~ \sum_{x\in b N}|\Cb_{B}(x)|\\
             &=\sum_{ b N\in{B/N}} ~ \sum_{x\in b N}|{(\Cb_{B}(x) N})/N||\Cb_N(x)|\\
             &\leqslant \sum_{b N\in{B/N}} ~ \sum_{x \in bN}|\Cb_{B/N}(x N)| |\Cb_N(x)|\\  
             &= \sum_{b N\in{B/N}}\Big(|\Cb_{B/N}(b N)|\sum_{x\in b N}|\Cb_N(x)|\Big)\\   
             &= \sum_{b N\in{B/N}}\Big(|\Cb_{B/N}(b N)|  \sum_{y\in N}|\Cb_{B}(y)\cap (b N)|\Big).\\                      
      \end{align*}
      Let $z \in   \Cb_B(y) \cap (b N)$.  Then $z N = b N$ as $z \in b  N$. Now, since $\Cb_{B}(y)$ is a subgroup of $(B, \circ)$ and $z \in \Cb_B(y)$, we get
      $$\Cb_{B}(y)\cap (z N) = (z \Cb_{B}(y)) \cap (z N) = z (\Cb_{B}(y)\cap N) = z   \Cb_N(y),$$
     which gives $|\Cb_{B}(y)\cap(b N)| = |\Cb_N(y)|$. By clubbing in these values in the preceding inequality, we get 
    \begin{align*}
        |B|^2\Pb(B) &\leqslant \sum_{bN\in B/N} \Big(|\Cb_{B/N}(b  N)| \sum_{y\in N}|\Cb_{N}(y)|\Big)\\
        &= \Big(\sum_{bN\in B/N} |\Cb_{B/N}(b  N)| \Big) \Big(\sum_{y\in N}|\Cb_{N}(y)|\Big),
    \end{align*}
      which, by the definition of the commuting probability,  can be written as
     $$|B|^2\Pb(B)\leqslant \big(|B/N|^2\Pb(B/N)\big)\big(|N|^2\Pb(N)\big).$$   
     Hence   $\Pb(B)\leqslant \Pb(N)\Pb(B/N)$ and the proof is complete. 
\end{proof}

\end{document}
